\section{Strong additional comparators and adaptation boundaries}\label{app:comparators}
The following controls belong to earlier independently audited experiments,
not the newly trained matched factorial model family. Keeping these identities
separate prevents a favorable acquisition or tuning result from being attributed
to the wrong training intervention.

\subsection{Ordinary-benefit and equal-allowance controls}
A robust/no-feedback price model with ordinary error-reduction acquisition was
evaluated on all six budgets and the original three seeds. The independently
checked 54-cell extension retained identical honest quotes, public scoring,
allocations and first-batch exact deviations. At 10\% capacity, reference
inflation was 0.01077, 0.01117 and 0.00778 on synthetic, JUCars and MUCars,
respectively. Reference-minus-SAV differences were $-0.00998$, $-0.00432$
and $0.00089$, with conditional shared-batch intervals
$[-0.01260,-0.00791]$, $[-0.00598,-0.00285]$ and $[0.00045,0.00130]$.
These intervals averaged the three fitted seeds before resampling batches;
they are not the ten-seed crossed intervals in the main factorial analysis.

An additional size/tuning study retained training caps 1,000, 2,500 and 5,000,
three fitting seeds, three training-internal gap-weight trials and twelve
residual epochs for both candidates. Its 972 parent tuning cells and 324
ordinary-benefit reference cells passed full audits. At the largest cap,
tuned SAV inflation was 20.41\% lower on MUCars but 30.09\% higher on
synthetic and 27.09\% higher on JUCars than the separately tuned ordinary-benefit
reference. The associated reference-minus-SAV intervals were
$[0.00097,0.00186]$, $[-0.00682,-0.00204]$ and $[-0.00333,-0.00125]$.
All scheduled caps and fixed/tuned roles remain in the archived evidence;
no favorable training size replaces the largest-cap comparison. This further
supports conditional gains rather than a consistent candidate advantage.

\subsection{Categorical support-margin imputation}
A whole-field adaptation of support-margin imputation (SMI)~\cite{smi} used
all allowed covariates and four optional-field groups. One deterministic
fitting generation per source trained a sixteen-model linear support-vector
regression family on complete-optional-field training rows. Scaling, category
vocabularies and target normalization were training-only. Missing categories
required a one-hot completion including an unknown token. If the reduced
sign-consistent model family or its optimization was unavailable, the method
used ordinary median/mode imputation. The original SMI guarantee is not
asserted on these fallback rows, and this is not an unchanged official reproduction.

An independent categorical-enumeration and scalar-interval auditor checked
94,527 report/domain cases and every final quotation and legal no-query report.
Training cohort sizes were 2,947, 2,674 and 3,368 for synthetic, JUCars and
MUCars. Across unrestricted, nonnegative and unit-box completion domains,
aggregate outcomes matched ordinary median/mode regression. Attacked fallback
fractions were 75.3\%, 87.3\% and 65.6\%, respectively. This diagnoses a limit
of the categorical adaptation's retained-family assumptions; it does not
refute the original SMI theorem. These no-query fits differ in cohort and
model class from the matched factorial predictor, so their higher inflation
cannot isolate a benefit of the present acquisition procedure.

\subsection{Acquisition-conditioned neighbor control}
ACO-1 adapted acquisition-conditioned expected-loss scoring~\cite{activeoracle}
to one-field regression queries and batch capacity. A training-only neighbor
bank estimated loss changes for sixteen observation patterns; unavailable
recorded bank values gave zero replay benefit. Distances used observed query
groups and permitted nonoptional features, with fixed category/numeric weights
and stable ties. It was a clean-room task adaptation, not an unchanged execution
of the original sequential classification method.

The independently audited nine-pair study contained 144 cells. The declared
main $k=25$ used every capacity; fixed $k=10$ and $k=50$ sensitivities used
10\% capacity only. No winning $k$ was selected. With the frozen robust
predictor and $k=25$, mean inflation was 0.01081 on synthetic, 0.01118 on
JUCars and 0.00716 on MUCars. Full-SAV-minus-ACO differences of 0.00995 and
0.00431 on synthetic and JUCars had conditional intervals
$[0.00725,0.01345]$ and $[0.00291,0.00588]$. On MUCars the point comparison
favored SAV by 0.00027, but the interval crossed zero. Thus a strong predictive
acquisition control can outperform the original candidate. These results
strengthen comparator coverage without implying that one-step ACO-1 dominates
the matched factorial policies or reproduces the original method's guarantee.
